\documentclass{article}
\usepackage{graphicx}
\usepackage[spanish,activeacute,es-tabla]{babel}
\usepackage[utf8x]{inputenc} % mac utf8
\usepackage{amsfonts}
\usepackage{minted}
\usepackage{ucs}
\usepackage{hyperref}


\begin{document}

\title{\textit{Teoría de la Información y Codificación}\\ \textbf{Práctica 3: Huffman como algoritmo de compresión}}
\author{José A. Montenegro Montes}

\maketitle


\section{Enunciado}

En esta práctica utilizaremos el algoritmo de Huffman para comprimir, basándonos en la codificación por símbolos.

Esta práctica será de utilidad para comparar posteriormente los algoritmos de compresión en bloques. 

La práctica está basada en el código del libro Algorithms 4 Edición(Autores Robert Sedgewick y Kevin Wayne), que puede ser consultado en \href{http://jabega.uma.es/search~S4*spi?/tAlgorithms/talgorithms/1%2C148%2C170%2CB/frameset&FF=talgorithms&4%2C%2C6#}{biblioteca}.



\section{Conclusiones}

Además de realizar la comparativa con los códigos en bloques, la práctica ofrece un ejemplo de lectura y escritura binaria y cómo es posible utilizar la librería diffutils.

El algoritmo es aplicado a los tres libros disponibles en el campus virtual: \textbf{Alice.txt}, \textbf{MobyDick.txt} y \textbf{Quixote.txt}. La práctica genera el archivo comprimido con el sufijo C.txt y el archivo descomprimido con el sufijo D.txt. Los sufijos puede ser modificados en el código. Por ejemplo, con la ejecución sobre el libro Alice.txt obtendremos Alice.txtC.txt como archivo comprimido y Alice.txtD.txt como archivo descomprimido.

Además de utilizar los libros anteriormente citados, sería interesante aplicar la práctica a un \textbf{documento aleatorio} y comparar su salida posteriormente con la codificación en bloques.

La salida del programa será la siguiente por defecto la siguiente:
\footnotesize
\begin{minted}[mathescape,
               linenos,
               numbersep=5pt,
               frame=lines,
               framesep=2mm]{csharp}
 
 
********************** 	 Alice.txt 	 **********************
Size regular: 152146 Size compress: 87816 Ratio: 0.5771825
Fichero Descomprimido correctamente!!! Size regular: 152146 Size decompress: 152146
Tiempo Comprimir: 65872000 nanosegundos. Tiempo Descomprimir: 56971000 nanosegundos.
****************************************************************
 
\end{minted}
\normalsize
donde nos ofrecen información del ratio de compresión, el tiempo utilizado en comprimir y descomprimir. Además verificamos que el archivo descomprimido es igual al original. 

Si ejecutamos la aplicación con la opción de estadísticas habilitada, nos mostrará la codificación utilizada, frecuencias y probabilidades de los caracteres del texto y finalmente la entropía y la longitud media obtenida.

\footnotesize
\begin{minted}[mathescape,
               linenos,
               numbersep=5pt,
               frame=lines,
               framesep=2mm]{csharp}
 
 
********************** 	 Alice.txt 	 **********************
Size regular: 152146 Size compress: 87816 Ratio: 0.5771825
Fichero Descomprimido correctamente!!! Size regular: 152146 Size decompress: 152146
Tiempo Comprimir: 65872000 nanosegundos. Tiempo Descomprimir: 56971000 nanosegundos.
****************************************************************


Caracter  (Frecuencia - Probabilidad) 	 Codificacion
------------------------------------------------------
10	 (3601 - 0.023668056)   01000
13	 (3601 - 0.023668056)   01001
26	 (1 - 6.572634E-6)   01011100110010010
32	 (28917 - 0.19006087)   00
33	 (450 - 0.0029576854)   111011011
34	 (113 - 7.4270763E-4)   11101101010
39	 (1763 - 0.011587554)   1110111
40	 (56 - 3.680675E-4)   111011001111
41	 (55 - 3.6149487E-4)   111011001110
42	 (60 - 3.9435804E-4)   01011100100
44	 (2418 - 0.015892629)   100111
45	 (669 - 0.0043970924)   11000010
46	 (978 - 0.0064280364)   0101011
48	 (1 - 6.572634E-6)   01011100110011000
51	 (1 - 6.572634E-6)   01011100110011001
58	 (233 - 0.0015314238)   010101000
59	 (194 - 0.0012750911)   1110110000
63	 (202 - 0.0013276722)   1110110010
65	 (638 - 0.0041933404)   10011011
66	 (91 - 5.9810973E-4)   11000011001
67	 (145 - 9.5303194E-4)   1001101001
68	 (192 - 0.0012619458)   1100001111
69	 (188 - 0.0012356553)   1100001110
70	 (74 - 4.8637492E-4)   10011010100
71	 (83 - 5.4552866E-4)   10011010111
72	 (284 - 0.0018666281)   010111011
73	 (733 - 0.004817741)   11000110
74	 (8 - 5.2581072E-5)   0101110011000
75	 (82 - 5.38956E-4)   10011010110
76	 (98 - 6.441182E-4)   11101100110
77	 (200 - 0.0013145269)   1110110001
78	 (120 - 7.887161E-4)   0101110000
79	 (176 - 0.0011567837)   1100001101
80	 (65 - 4.272212E-4)   01011100101
81	 (84 - 5.521013E-4)   11000011000
82	 (140 - 9.2016876E-4)   0101110101
83	 (218 - 0.0014328342)   1110110100
84	 (472 - 0.0031022832)   01010101
85	 (66 - 4.3379384E-4)   01011100111
86	 (42 - 2.7605065E-4)   010111001101
87	 (237 - 0.0015577143)   010101001
88	 (4 - 2.6290536E-5)   010111001100101
89	 (114 - 7.4928027E-4)   11101101011
90	 (1 - 6.572634E-6)   01011100110010011
91	 (2 - 1.3145268E-5)   0101110011001000
93	 (2 - 1.3145268E-5)   0101110011001101
95	 (4 - 2.6290536E-5)   010111001100111
96	 (1109 - 0.007289051)   1001100
97	 (8153 - 0.053586688)   0111
98	 (1384 - 0.009096526)   1100000
99	 (2254 - 0.014814718)   100001
100	 (4739 - 0.031147713)   10010
101	 (13386 - 0.08798128)   1101
102	 (1927 - 0.012665466)   010110
103	 (2448 - 0.016089808)   101001
104	 (7090 - 0.046599977)   11111
105	 (6781 - 0.044569034)   11100
106	 (138 - 9.070235E-4)   0101110100
107	 (1076 - 0.0070721544)   0101111
108	 (4618 - 0.030352425)   10001
109	 (1907 - 0.012534013)   010100
110	 (6896 - 0.045324884)   11110
111	 (7970 - 0.052383896)   0110
112	 (1459 - 0.009589474)   1100010
113	 (125 - 8.2157925E-4)   0101110001
114	 (5297 - 0.034815244)   10101
115	 (6282 - 0.04128929)   11001
116	 (10217 - 0.067152604)   1011
117	 (3402 - 0.022360101)   111010
118	 (804 - 0.005284398)   11000111
119	 (2438 - 0.016024083)   101000
120	 (144 - 9.464593E-4)   1001101000
121	 (2149 - 0.014124591)   100000
122	 (77 - 5.0609285E-4)   10011010101

Entropia:       4.5675883
Longitud Media: 4.6123915
 
\end{minted}



\section{Código}


\subsection*{Clase FuenteFile.} 

\inputminted[mathescape,
               numbersep=5pt,
               frame=lines,
               framesep=2mm]{java}{practica3/FuenteFile.java}

               %%%%%%%%%%%%%%%%%%%%%%%
 
\subsection*{Clase Huffman.} 
\inputminted[mathescape,
               numbersep=5pt,
               frame=lines,
               framesep=2mm]{java}{practica3/Huffman.java}
               
%%%%%%
\subsection*{Clase MinPQ.} 
\inputminted[mathescape,
               numbersep=5pt,
               frame=lines,
               framesep=2mm]{java}{practica3/MinPQ.java}
               
%%%%%%
 
\subsection*{Clase BinaryIn.} 
\inputminted[mathescape,
               numbersep=5pt,
               frame=lines,
               framesep=2mm]{java}{practica3/BinaryIn.java}
               
%%%%%%
\subsection*{Clase BinaryOut.} 
\inputminted[mathescape,
               numbersep=5pt,
               frame=lines,
               framesep=2mm]{java}{practica3/BinaryOut.java}
               
%%%%%%

\end{document}
